% !PRINTED 529
\weberpage{529}
%<<CONT
and thus obtain, since by (10)
\weq{15}{\varphi_+(0) = \varphi_-(0) = \varphi(0) = \int_0^\infty e^{\pi iyv(\omega_1+\omega_2)}\left(\frac{v}{2}\right)^{2s-2}dv}
holds, from (11), (12) and (14):
\weq{16}{\Sigma_{-\infty,\infty}^{x}W = \varphi(0) + \Sigma_{1,\infty}^{\nu}[\varphi_+(\nu)+\varphi_-(\nu)].}

If we introduce this into (8), then $M_s$ splits into two parts:
\weq{17}{M_s = P_s + Q_s,}
if
\weq{18}{P_s = \frac{(2\pi)^{2s}}{A^s\Gamma(s)\Gamma(s)}\Sigma_{1,\infty}^{y}\varphi(0),}
\weq{19}{Q_s = \frac{(2\pi)^{2s}}{A^s\Gamma(s)\Gamma(s)}\Sigma_{1,\infty}^{y}\Sigma_{1,\infty}^{\nu}[\varphi_+(\nu)+\varphi_-(\nu)]}
are set.

By (15), if for $v$ we introduce a new variable of integration
$t$ by the equation
\[ \pi iv(\omega_1+\omega_2) = -\frac{2\pi v\sqrt{m}}{A} = -t \]
and carry out the summation over $y$, we have:
\[ \Sigma_{1,\infty}^{y}\varphi(0) = \frac{2A^{2s-1}}{(4\pi\sqrt{m})^{2s-1}}\int_0^\infty\frac{t^{2s-2}\,dt}{e^t-1}, \]
and therefore:
\weq{20}{P_s = \frac{(2\pi)^{2s}A^{s-1}}{(4\pi\sqrt{m})^{2s-1}\Gamma(s)^2}\,2\int_0^\infty\frac{t^{2s-2}\,dt}{e^t-1}.}

From this the limiting value $P$ can easily be determined. For
the integral
\weq{21}{\int_0^\infty t^{2s-2}e^{-t}\left(\frac{1}{1-e^{-t}}-\frac{1}{t}\right)dt,}
is continuous up to $s=1$, since the expression standing within the
parenthesis retains a finite value for $t=0$ and $t=\infty$, and it
therefore has the limiting value:
\[ \int_0^\infty e^{-t}\left(\frac{1}{1-e^{-t}}-\frac{1}{t}\right)dt = -\Gamma'(1) = 0{,}5772\ldots \]
